\documentclass[tikz,border=8pt]{standalone}
\usepackage{fontspec}
\usepackage{xeCJK}
\IfFontExistsTF{Segoe UI}{\setmainfont{Segoe UI}}{\setmainfont{Arial}}
\IfFontExistsTF{Microsoft JhengHei}{\setCJKmainfont{Microsoft JhengHei}}{\IfFontExistsTF{Noto Sans CJK TC}{\setCJKmainfont{Noto Sans CJK TC}}{\setCJKmainfont{Noto Sans TC}}}
\IfFontExistsTF{Microsoft JhengHei}{\setCJKsansfont{Microsoft JhengHei}}{\IfFontExistsTF{Noto Sans CJK TC}{\setCJKsansfont{Noto Sans CJK TC}}{\setCJKsansfont{Noto Sans TC}}}
\xeCJKsetup{CJKglue={\hskip 0.045em plus 0.015em minus 0.005em}, CJKecglue={\hskip 0.08em plus 0.02em}}
\IfFileExists{../../../FontAwesome.otf}{%
  \newfontfamily\FA[Path=../../../]{FontAwesome.otf}%
}{%
  \newfontfamily\FA{Segoe UI Symbol}%
}
\newcommand{\faIcon}[1]{{\FA\symbol{#1}}}
\usetikzlibrary{arrows.meta,positioning}
\definecolor{paper}{HTML}{FFFDF8}
\definecolor{navy}{HTML}{0F172A}
\definecolor{muted}{HTML}{334155}
\definecolor{green}{HTML}{00856A}
\definecolor{greensoft}{HTML}{DCFCE7}
\definecolor{gold}{HTML}{D97706}
\definecolor{goldsoft}{HTML}{FEF3C7}
\definecolor{line}{HTML}{CBD5E1}
\begin{document}
\begin{tikzpicture}[x=1cm,y=1cm,>=Latex]
\path[fill=paper] (0,0) rectangle (16,9.6);
\node[font=\bfseries\fontsize{18}{22}\selectfont,text=navy,align=center,text width=14.6cm] at (8,8.82) {會計資料卡};
\node[font=\fontsize{10.5}{13}\selectfont,text=muted,align=center,text width=13.8cm] at (8,8.12) {不是交程式碼，而是交出可以被信任的資料說明};

\draw[rounded corners=12pt,fill=white,draw=line,line width=.9pt] (3.0,1.95) rectangle (13.0,7.35);
\draw[rounded corners=12pt,fill=greensoft,draw=green,line width=.9pt] (3.0,6.45) rectangle (13.0,7.35);
\node[font=\bfseries\fontsize{13}{16}\selectfont,text=navy,align=center,text width=8.6cm] at (8,6.9) {一張資料表，至少要能回答六件事};

\foreach \x/\y/\title/\body in {
4.55/5.55/來源/從哪個網站或檔案取得,
8.00/5.55/單位/元、千元、百分比是否一致,
11.45/5.55/期間/月、季、年能否對齊,
4.55/3.88/例外/缺值與跳號怎麼處理,
8.00/3.88/核對/抽查哪些列與合計,
11.45/3.88/解讀/資料回答哪個會計問題}{
  \draw[rounded corners=7pt,fill=goldsoft,draw=gold,line width=.9pt] (\x-1.28,\y-.68) rectangle (\x+1.28,\y+.68);
  \node[font=\bfseries\fontsize{9.2}{11}\selectfont,text=navy,align=center,text width=2.12cm] at (\x,\y+.32) {\title};
  \node[font=\fontsize{7.1}{8.6}\selectfont,text=muted,align=center,text width=2.12cm] at (\x,\y-.25) {\body};
}

\draw[->,line width=1pt,draw=navy] (4.55,5.0)--(4.55,4.5);
\draw[->,line width=1pt,draw=navy] (8.0,5.0)--(8.0,4.5);
\draw[->,line width=1pt,draw=navy] (11.45,5.0)--(11.45,4.5);
\draw[draw=line,line width=.9pt] (3.45,2.7)--(12.55,2.7);
\node[font=\fontsize{10.5}{13}\selectfont,text=navy,align=center,text width=13.4cm] at (8,2.28) {資料卡讓學生把「有用」和「能用」分開};
\end{tikzpicture}
\end{document}
